% Generated draft: requires mathematical and source review.
\documentclass[11pt]{article}
\usepackage[margin=25mm]{geometry}
\usepackage{fontspec}
\setmainfont{Latin Modern Roman}
\usepackage{amsmath,amssymb,mathtools,unicode-math}
\setmathfont{Latin Modern Math}
\usepackage{enumitem,needspace,xurl,hyperref}
\setlength{\parindent}{0pt}
\setlength{\parskip}{5pt}
\setlength{\emergencystretch}{3em}
\newcommand{\N}{\mathbb{N}}
\newcommand{\Z}{\mathbb{Z}}
\newcommand{\Q}{\mathbb{Q}}
\newcommand{\R}{\mathbb{R}}
\newcommand{\C}{\mathbb{C}}
\newcommand{\F}{\mathbb{F}}
\newcommand{\A}{\mathbb{A}}
\newcommand{\PP}{\mathbb{P}}
\newcommand{\E}{\mathbb{E}}
\newcommand{\eps}{\varepsilon}
\newcommand{\dd}{\,\mathrm{d}}
\begin{document}
\section*{UnsolvedMath ID: 30002328}
\textbf{Source problem code:} catalogue entry
\subsection{Definitions and mathematical statement}
The underlying spatial lattice is \(\Z^2\). A branching random walk consists of particles that move on this lattice and reproduce according to specified probabilistic rules. Here a cutoff additionally limits the number of particles retained at each site. The source does not specify these rules, the cutoff value or mechanism, the time parameter, or the initial configuration.

Once such a process is specified, let \(\eta_t(x)\in\Z_{\geq 0}\) denote the number of retained particles at site \(x\in\Z^2\) at time \(t\). If ``occupied set'' means the sites currently containing particles, its definition would be
\[
\mathcal O_t^{\mathrm{inst}}=\{x\in\Z^2:\eta_t(x)>0\}.
\]
This is a conditional definition: the supplied record does not distinguish current occupancy from occupation up to time \(t\), nor does it specify the boundary convention.

\emph{Source question.} For a branching random walk on \(\Z^2\) with a cutoff limiting the number of particles at each site, is the boundary of the occupied set in the Kardar--Parisi--Zhang (KPZ) universality class?

\emph{Unresolved mathematical qualifications.} The supplied record does not define membership in the KPZ universality class for this boundary. In particular, it gives no boundary observable, centering, scaling, asymptotic regime, or limiting assertion. It also does not specify whether the question concerns one particular model or every model satisfying additional hypotheses. Consequently, this record determines a qualitative conjectural question but does not determine a unique precise probabilistic statement.
\subsection{Short English statement}
The question asks whether the growing occupied boundary of a branching random walk on \(\Z^2\), subject to a per-site particle cutoff, exhibits KPZ behavior. The model, occupied-set convention, and mathematical meaning of KPZ behavior require specification before a precise conjecture can be stated.
\subsection{Sources}
\begin{itemize}
\item Supplied local UnsolvedMath problem registry record, permanent ID \(30002328\), exact source problem code \texttt{catalogue entry}, title ``KPZ Boundaries in Cutoff Branching Random Walk,'' in the Oberwolfach Reports Open Problems collection. No individual proposer is supplied; the proposed year is recorded as 2013.
\item Supplied source citation: \emph{Extremes in Branching Random Walk and Branching Brownian Motion} (2014), Oberwolfach Reports, pp.~1205--1251, DOI \texttt{10.4171/owr/2013/20}. \url{https://doi.org/10.4171/owr/2013/20}. The background describes these reports as records of mathematical workshops and their open-problem discussions.
\item Supplied literature source: ``Oberwolfach Report 20/2013,'' \url{https://doi.org/10.4171/owr/2013/20}. The literature assessment dated \texttt{2026-08-21} labels the problem open and describes the cutoff branching-random-walk KPZ claim as an explicitly stated conjecture, with no proof or disproof found by the reported search for the model as formulated. This is supplied dated curation derived from the OWR record, not an independently verified literature assessment; time-sensitive status and cited claims require independent verification.
\item The supplied statement status is \texttt{corrected\_verified}. Its verification note reports comparison with the cited official Oberwolfach report and normalization of notation, context, and question form without changing the source-level claim. This is supplied verification metadata, not a verification performed for this draft.
\end{itemize}
\end{document}
